\documentclass[tikz,border=10pt]{standalone}
\usepackage{geometricfigures}
\usetikzlibrary{positioning, mindmap}

% The fields of `NonlinearSolver`, in the order of its definition in SimpleSolvers
% (src/nonlinear/nonlinear_solver.jl), each with the bound of its type parameter.
\begin{document}
\begin{tikzpicture}[mindmap, concept color=fg, text=bg,%
                   level 1 concept/.append style={sibling angle=45},%
                   level 2 concept/.append style={sibling angle=40}]
  \path[mindmap,concept color=fg,text=bg]
    node[concept, text=bg] {\texttt{NonlinearSolver <: \\AbstractSolver}}
    [clockwise from=90]
    % note that `sibling angle' can only be defined in
    % `level 1 concept/.append style={}'
    % note that the `concept color' is passed to the `child'(!)
    child[concept color=mblue] {node[concept] {\texttt{nonlinear-\\problem::\\Nonlinear-\\Problem}} [clockwise from=110] child { node[concept] {\texttt{F::\\Callable}}} child { node[concept] {\texttt{J::\\Callable\\or Missing}}}}
    child[concept color=mblue!50!mgreen] {node[concept] {\texttt{linear-\\problem::\\Abstract-\\Linear-\\Problem}} }
    child[concept color=mgreen] {node[concept] {\texttt{jacobian::\\Jacobian}} }
    child[concept color=mgreen!50!morange] {node[concept] {\texttt{linear-\\solver::\\Abstract-\\Linear-\\Solver}} }
    child[concept color=morange] {node[concept] {\texttt{linesearch::\\Linesearch}} }
    child[concept color=mred] {node[concept] {\texttt{method::\\Nonlinear-\\Solver-\\Method}} }
    child[concept color=mred!50!mpurple] {node[concept] {\texttt{cache::\\Abstract-\\Nonlinear-\\Solver-\\Cache}} }
    child[concept color=mpurple] {node[concept] {\texttt{config::\\Options}}};
\end{tikzpicture}
\end{document}
